\documentclass[11pt,a4paper]{article}

% --- Paket tanımları ---
\usepackage[utf8]{inputenc}
\usepackage[T1]{fontenc}
\usepackage[turkish]{babel}
\usepackage{amsmath, amssymb}
\usepackage{graphicx}
\usepackage{hyperref}
\usepackage{xcolor}
\usepackage{geometry}
\usepackage{listings}
\usepackage{caption}
\usepackage{booktabs}
\usepackage{enumitem}
\usepackage{float}

\geometry{margin=2.5cm}
\hypersetup{colorlinks=true, linkcolor=blue, urlcolor=blue, citecolor=blue}

% --- Python listing stili ---
\definecolor{codegreen}{rgb}{0,0.6,0}
\definecolor{codegray}{rgb}{0.5,0.5,0.5}
\definecolor{codepurple}{rgb}{0.58,0,0.82}
\definecolor{backcolor}{rgb}{0.96,0.96,0.96}

\lstdefinestyle{pythoncode}{
  backgroundcolor=\color{backcolor},
  commentstyle=\color{codegreen},
  keywordstyle=\color{magenta},
  numberstyle=\tiny\color{codegray},
  stringstyle=\color{codepurple},
  basicstyle=\ttfamily\footnotesize,
  breakatwhitespace=false,
  breaklines=true,
  captionpos=b,
  keepspaces=true,
  numbers=left,
  numbersep=5pt,
  showspaces=false,
  showstringspaces=false,
  showtabs=false,
  tabsize=2,
  language=Python,
  extendedchars=true,
  inputencoding=utf8,
  literate=
    {ç}{{\c{c}}}1 {ş}{{\c{s}}}1 {ğ}{{\u{g}}}1 {ı}{{\i}}1 {ö}{{\"o}}1 {ü}{{\"u}}1
    {Ç}{{\c{C}}}1 {Ş}{{\c{S}}}1 {Ğ}{{\u{G}}}1 {İ}{{\.{I}}}1 {Ö}{{\"O}}1 {Ü}{{\"U}}1
    {σ}{{$\sigma$}}1 {μ}{{$\mu$}}1 {π}{{$\pi$}}1 {α}{{$\alpha$}}1 {β}{{$\beta$}}1
    {ω}{{$\omega$}}1 {Γ}{{$\Gamma$}}1 {Δ}{{$\Delta$}}1 {θ}{{$\theta$}}1 {φ}{{$\varphi$}}1
    {≈}{{$\approx$}}1 {±}{{$\pm$}}1 {→}{{$\rightarrow$}}1 {²}{{$^2$}}1
}
\lstset{style=pythoncode}

% --- Kapak ---
\title{\textbf{Mikro-Doppler İmzaları Üzerinde Model-Tabanlı Veri Artırma Yönteminin Gerçeklenmesi}\\[0.5em]
\large Rojhani vd. (IEEE TMTT 2023) Makalesinin Python ile Sentetik FMCW Radar Pipeline Reproduksiyonu}
\author{Bahadır Çeliktaş\\
TEL503 — Yüksek Lisans Projesi}
\date{\today}

\begin{document}
\maketitle
\tableofcontents
\newpage

% =============================================================
\section{Özet}
% =============================================================
Bu rapor, Rojhani \emph{vd.}'in IEEE Transactions on Microwave Theory and Techniques (Mayıs 2023) sayısında yayımlanan
\emph{``Model-Based Data Augmentation Applied to Deep Learning Networks for Classification of Micro-Doppler Signatures Using FMCW Radar''}
başlıklı makalesinin sentetik veri üretim zinciri ve sınıflandırma deneylerinin Python dilinde gerçeklenmesini detaylı biçimde belgeler.

Makale, küçük insansız hava araçlarının (UAV) FMCW (Frequency Modulated Continuous Wave) radar
imzalarını fiziksel/deterministik bir model ile sentezleyerek derin öğrenme tabanlı sınıflandırıcılarda
veri artırma (data augmentation) yöntemi olarak kullanmayı önermektedir. Bu projede:
\begin{itemize}
    \item Makaledeki tüm temel denklemler (Eq.~6, 10, 11, 12, 14, 15, 16, 19--22) ve
          parametre tabloları (Table I, Table II) Python NumPy üzerinde birebir kodlanmıştır.
    \item Önerilen CNN mimarisi (Fig.~7) PyTorch ile aynı katman yapısında inşa edilmiştir.
    \item Makalede karşılaştırma için kullanılan üç klasik makine öğrenmesi sınıflandırıcısı
          (SVM, Naive Bayes, XGBoost) scikit-learn ve XGBoost kütüphaneleri ile kurulmuştur.
    \item Makalenin gerçek ölçüm veri seti kamuya açık olmadığından,
          ölçüm setinin yerine sentetik bir \emph{ölçüm proxy'si} (sabit $\sigma_P=0{,}20$ + 20\,dB AWGN)
          tanımlanmış ve cross-domain test bu set üzerinde yapılmıştır.
    \item Makalenin temel deneyleri (Fig.~11 sentez/ölçüm karşılaştırması, Fig.~13 $\sigma_P$ tarama eğrisi,
          Table~IV ML sınıflandırıcı tablosu, Fig.~14 confusion matrix) ya birebir ya da yapısal olarak
          uyumlu biçimde elde edilmiştir.
\end{itemize}

% =============================================================
\section{Giriş ve Projenin Amacı}
% =============================================================
\subsection{Problem Tanımı}
Mini insansız hava araçları (mini-UAV; helikopter, quadcopter vb.) düşük uçuş yüksekliği, küçük radar
kesit alanı (RCS) ve düşük hız nedeniyle radar tabanlı algılama sistemleri için zorlu hedeflerdir.
Bu hedefleri ayırt etmek için yaygın olarak \emph{mikro-Doppler imza} (micro-Doppler signature, $\mu$D)
kullanılır: hedefin gövde Doppler bileşeni dışında, dönen pervane/rotor gibi mekanik parçalardan
gelen ek frekans bileşenleridir.

Derin öğrenme tabanlı sınıflandırıcılar (özellikle CNN) bu görevlerde yüksek başarı
gösterse de \textbf{büyük miktarda etiketli ölçüm verisine ihtiyaç duyarlar}.
Gerçek radar ölçüm kampanyaları pahalı, yavaş ve hava şartları/test alanı kısıtlarına bağlıdır.
Rojhani \emph{vd.} bu probleme \textbf{model tabanlı (deterministik) veri artırma}
yaklaşımını önerir: hedefin fiziksel/mekanik özelliklerini içeren bir matematiksel modelden
sentetik dataset üretmek ve CNN'i bu sentetik dataset ile eğitmek.

\subsection{Projenin Hedefleri (Başarı Ölçütleri)}
\begin{description}
    \item[Ölçüt 1 — Sinyal modelinin fiziksel tutarlılığı:] Üretilen FMCW IF sinyali, hem ana
          Doppler bantını hem de pervane dönüş frekanslarındaki mikro-Doppler yan bantlarını
          doğru frekanslarda içermelidir.
    \item[Ölçüt 2 — Hedef tipi ayırt edilebilirliği:] Tek motorlu (HELI) ve quadcopter (QUAD) sınıfları
          arasındaki imza farkı sınıflandırıcılar tarafından öğrenilebilir olmalıdır.
    \item[Ölçüt 3 — Literatür ile nicel uyum:] CNN sınıflandırma doğruluğu, makalenin raporladığı
          \%78{,}68 değerine yaklaşmalıdır.
\end{description}

% =============================================================
\section{Telekom Domain Bilgileri (Makaleden)}
% =============================================================
\subsection{FMCW Radar Çalışma Prensibi}
FMCW radar, sürekli olarak frekansı zamanla doğrusal değişen bir \emph{chirp} sinyali yayınlayan
ve geri yansıyan sinyali iletilen sinyal ile karıştırarak (mixer) düşük frekanslı bir
\textbf{IF (Intermediate Frequency)} sinyali üreten bir radar mimarisidir. IF sinyalinin frekansı,
hedefin radara olan mesafesi ile orantılıdır:
\begin{equation}
R = \frac{c_0 \cdot f_{\text{IF}}}{2\mu}
\label{eq:range}
\end{equation}
Burada $c_0$ ışık hızı, $\mu$ chirp frekans rampa katsayısı (Hz/s), $f_{\text{IF}}$ ise
IF sinyalindeki tepe frekansıdır. Eq.~\ref{eq:range} makalenin Eq.~(14)'üdür.

\subsection{Mikro-Doppler İmza}
Hedef tamamen rijit bir cisim ise IF spektrumda yalnızca tek bir tepe gözlenir.
Dönen rotor gibi mekanik bir parça varsa, bu parçadan saçılan elektromanyetik dalga
zamanla değişen bir mesafeye sahiptir; sonuç olarak ek frekans bileşenleri (yan bantlar)
oluşur. Bu yan bantlar topluca \emph{mikro-Doppler imzası} ($\mu$D) olarak adlandırılır.
Pervane/rotor frekansı ve geometrisine özgü olduğundan farklı UAV sınıflarını
ayırt etmede kullanılabilir.

\subsection{IF Sinyalinin Matematiksel Modeli (Eq. 6)}
Makale, boresight'ta (radarın bakış doğrultusunda) duran bir UAV için IF sinyalini şu şekilde verir:
\begin{equation}
s_{\text{IF}}(t) = \sum_{k=1}^{N_{\text{eng}}} \Gamma_k \cos\!\left[\frac{8\pi\mu}{c_0}\left(R_0 + \sum_{j=1}^{4} A_{kj} \cos(\omega_{kj} t + \phi_{kj})\right) t\right]
\label{eq:eq6}
\end{equation}
Bu denklemde:
\begin{itemize}
    \item $N_{\text{eng}}$: rotor (motor) sayısı (HELI için 1, QUAD için 4).
    \item $\Gamma_k$: $k$. rotorun beklenen RCS (Radar Kesit Alanı).
    \item $R_0$: hedefin sabit (gövde) menzili (m).
    \item $A_{kj}, \omega_{kj}, \phi_{kj}$: $k$. rotorun $j$. vibrasyonel bileşeninin
          genlik, frekans ve faz değerleri.
    \item $j$ indeksi sırasıyla şu bileşenlere karşılık gelir: $\alpha$ (roll), $\beta$ (pitch),
          $|\omega_\alpha + \omega_\beta|$ (toplam karışım), $|\omega_\alpha - \omega_\beta|$ (fark karışım).
\end{itemize}
Yani her motor için 13 skalar (1 RCS $\Gamma_k$ + 4 bileşen $\times$ 3 parametre) tanımlıdır.

\subsection{Vibrasyonel Bileşenlerin Açılımı (Eq. 10)}
Tek bir motorun vibrasyonel parametre vektörü (Eq.~10):
\begin{equation}
\bar{x}_k = \langle \Gamma_k,\; A_{k\alpha}, \omega_{k\alpha}, \phi_{k\alpha},\;
                              A_{k\beta}, \omega_{k\beta}, \phi_{k\beta},\;
                              A_{k(\alpha+\beta)}, |\omega_{k\alpha}+\omega_{k\beta}|, \phi_{k(\alpha+\beta)},\;
                              A_{k(\alpha-\beta)}, |\omega_{k\alpha}-\omega_{k\beta}|, \phi_{k(\alpha-\beta)} \rangle^T
\label{eq:eq10}
\end{equation}

Burada karışım terimlerinin frekans katsayıları (Eq.~11):
\begin{equation}
C_{\omega_{k(\alpha+\beta)}} = \frac{\omega_{k(\alpha+\beta)}}{|\omega_{k\alpha}+\omega_{k\beta}|}, \qquad
C_{\omega_{k(\alpha-\beta)}} = \frac{\omega_{k(\alpha-\beta)}}{|\omega_{k\alpha}-\omega_{k\beta}|}
\end{equation}
şeklinde tanımlanır. Bu katsayılar Table~II'de \%100, \%99, \%92, \%104 gibi yüzdelik değerlerle
verilmiştir.

\subsection{Veri Artırma için Parametre Perturbasyonu (Eq. 12)}
Sınıf-içi varyasyonu modellemek için her parametreye, kendi büyüklüğüyle orantılı standart sapmaya
sahip Gauss gürültüsü eklenir:
\begin{equation}
\bar{n}_{i,P} = \big(\,\mathcal{N}(0, \sigma_P |x|_1),\,
                    \mathcal{N}(0, \sigma_P |x|_2),\, \ldots\,\big)^T,
                    \quad \sigma_P \in [0,1]
\label{eq:eq12}
\end{equation}
Burada $\sigma_P$ \emph{tek skaler}, ne kadar artırma yapılacağını kontrol eden hiperparametredir.
$\sigma_P=0$ deterministik (tek bir profil), $\sigma_P=0{,}60$ ise çok geniş bir dağılım üretir.
Makale Section IV-C-1'de $\sigma_P=0{,}20$ değerinin ölçüm setine en iyi uyduğunu deneysel olarak
göstermektedir.

\subsection{MTI (Moving Target Indicator) Filtresi (Eq. 15)}
Gerçek ölçümlerde hareketsiz hedefler (zemin yansıması, statik gövde) ve düşük hızlı clutter
istenmeyen bileşenlerdir. MTI filtresi bunları bastırmak için ardışık radar görüntüleri arasında
fark alır:
\begin{equation}
I_{\text{MTI}}(r,\theta) = I(r,\theta) - \frac{1}{N_{\text{MTI}}}\sum_{i=1}^{N_{\text{MTI}}} L_{-i}(r,\theta)
\label{eq:eq15}
\end{equation}
Makalede $N_{\text{MTI}} = 10$ frame kullanılır (Table~I).

\subsection{Açısal Entegrasyon (Eq. 16)}
MIMO array, açıya bağlı (range-angle) iki boyutlu bir görüntü üretir. Boresight hedefine ait 1-D menzil
profili elde etmek için açısal pencere boyunca (paperde $\Delta\theta = 7^\circ$) ortalama alınır:
\begin{equation}
\bar{s}(r) = \frac{1}{\Delta\theta}\int_{\theta - \Delta\theta/2}^{\theta + \Delta\theta/2} I_{\text{MTI}}(r, \theta')\, d\theta'
\end{equation}

\subsection{Donanım Parametreleri (Table I)}
\begin{table}[H]
\centering
\caption{Makaledeki FMCW MIMO radar platform parametreleri (Table I).}
\label{tab:tableI}
\begin{tabular}{ll}
\toprule
\textbf{Parametre} & \textbf{Değer} \\
\midrule
Ramp frekans aralığı & 75 -- 79 GHz (taşıyıcı $f_c \approx 77$ GHz) \\
RF çıkış gücü & 10 dBm \\
RX anten sayısı & 8 \\
TX anten sayısı & 4 \\
Ramp süresi (yukarı) & 2{,}56 ms \\
Ramp süresi (aşağı) & 0{,}32 ms \\
Ramp katsayısı $\mu$ & $1{,}56 \times 10^{12}$ Hz/s \\
Frame başına ramp & 1000 \\
Chirp tekrar aralığı (CRI) & 200 ms \\
Chirp başına örnek $N_{\text{samples}}$ & 504 \\
IF örnekleme frekansı $F_s$ & 200 kHz \\
FFT noktası sayısı $N_{\text{FFT}}$ & 4096 \\
FFT işlem kazancı & 36{,}12 dB \\
MTI frame derinliği $N_{\text{MTI}}$ & 10 \\
\bottomrule
\end{tabular}
\end{table}

\subsection{Hedef Sınıfları ve Model Parametreleri (Table II)}
Makalenin Table~II'si iki referans hedefe ait parametre vektörlerini verir:
\begin{itemize}
    \item \textbf{HELI} (RadioFLY plastik helikopter, 1 motor): $\Gamma_2 = \Gamma_3 = \Gamma_4 = 0$,
          yalnızca Engine~1 aktif. $\Gamma_1 = 0{,}044$, $\omega_{1\alpha} = 2{,}69 \times 10^3$ rad/s ($\sim 428$ Hz),
          $\omega_{1\beta} = 3{,}98 \times 10^3$ rad/s ($\sim 633$ Hz).
    \item \textbf{QUAD} (DJI Mavic Mini 2, 4 motor): Dört motorun tümü aktif. $\Gamma_k$ değerleri
          $0{,}001 - 0{,}011$ aralığında, $\omega$ değerleri 1{,}77\,kHz -- 5{,}07\,kHz aralığında.
\end{itemize}
Bu parametreler doğrudan kodumuzda \texttt{HELI\_REF\_PARAMS} ve \texttt{QUAD\_REF\_PARAMS}
sözlüklerine birebir kopyalanmıştır.

\subsection{Veri Üretim Zincirleri}
Makalede iki farklı zincir tanımlanır:
\begin{description}
    \item[Fig. 2 — Sentetik zincir:] Eq.~6 ile $s_{\text{IF}}(t)$ üret $\rightarrow$ FFT $\rightarrow$
          tepe etrafından 400 örneklik menzil profili al. Bu zincir CNN/ML eğitim verisi içindir.
    \item[Fig. 6 — Ölçüm zinciri:] Çoklu frame topla $\rightarrow$ 2-D FFT $\rightarrow$ MTI
          (Eq.~15) $\rightarrow$ açısal entegrasyon (Eq.~16) $\rightarrow$ tepe etrafından
          400 örnek. Bu zincir gerçek ölçümün işlenmesine karşılık gelir.
\end{description}

\subsection{Önerilen CNN Mimarisi (Fig. 7)}
\begin{itemize}
    \item Giriş: $400 \times 1$ menzil profili $\rightarrow$ $1 \times 20 \times 20$'ye yeniden şekillendir.
    \item \textbf{Conv1}: $K=3$, 32 filtre, ReLU $\rightarrow$ çıktı $32 \times 18 \times 18$.
    \item \textbf{MaxPool1}: $2 \times 2$ $\rightarrow$ $32 \times 9 \times 9$.
    \item \textbf{Conv2}: $K=3$, 64 filtre, ReLU $\rightarrow$ $64 \times 7 \times 7$.
    \item \textbf{AvgPool2}: $7 \times 7$ $\rightarrow$ $64 \times 1 \times 1$.
    \item \textbf{Flatten} $\rightarrow$ $64$ boyutlu vektör.
    \item \textbf{FC1}: $64 \rightarrow 32$, ReLU.
    \item \textbf{FC2}: $32 \rightarrow 2$ (sınıf logitleri).
    \item Softmax + Cross-Entropy kayıp (Eq.~19, 20).
    \item Adam optimizer, $\text{lr} = 10^{-3}$, 50 epoch (Section IV-C-2).
\end{itemize}

\subsection{Performans Metrikleri (Eq. 21--22)}
\begin{equation}
\text{acc} = \frac{TP + TN}{TP + FP + FN + TN}, \qquad
\text{prec} = \frac{TP}{TP + FP}
\end{equation}

\subsection{Makalenin Raporladığı Sonuçlar}
\begin{itemize}
    \item Table~IV: SVM, Naive Bayes, XGBoost için $\sigma_P \in \{0, 0{,}05, 0{,}10, 0{,}20, 0{,}30\}$ tarama sonuçları.
          ``New'' (136 örnek) ölçüm setinde optimum nokta $\sigma_P = 0{,}20$.
    \item Fig.~13: $\sigma_P$ vs accuracy eğrisi, üç sınıflandırıcı için ölçüm seti üzerinde
          $\sigma_P \approx 0{,}20$'de tepe yapan çan eğrisi.
    \item CNN ile $\sigma_P=0{,}20$ ve 16k/sınıf eğitim seti (toplam 32000 örnek) ile gerçek
          ölçümler üzerinde \%78{,}68 doğruluk (Fig.~14).
    \item Conventional augmenter (\emph{tsaug}) ile sadece \%66{,}18 (Fig.~15), ve dengesiz
          (HELI hassasiyeti yalnızca \%15{,}55).
\end{itemize}

% =============================================================
\section{Bizim Telekom Tasarım Kararlarımız ve Yorumlarımız}
% =============================================================
Makalede bazı noktalar açık değildir veya farklı yorumlara açıktır. Bu bölüm, kodda yapılan
kritik telekom kararlarını ve nedenlerini belgeler.

\subsection{Eq. 6'daki $8\pi\mu/c_0$ Katsayısı}
Eq.~6'da yazılı $8\pi\mu/c_0$ katsayısı $4\pi\mu/c_0$ olarak alındı.
\textbf{Gerekçe}: Eq.~14 ($R = c_0 \cdot f / 2\mu$) ile tutarlılık için chirp slope üzerinden
elde edilen IF frekansı $f_{\text{IF}} = (2\mu/c_0)\cdot R$ olmalıdır; cosinüsün argümanı
$2\pi f_{\text{IF}} t = (4\pi\mu/c_0)\cdot R\cdot t$ verir. ``8'' katsayısı ile peak menzili
2$\times$ kayar. Kodda:
\begin{equation*}
K = \frac{4\pi\mu}{c_0}, \qquad \text{phase}(t) = K\cdot(R_0 + \text{vib}(t))\cdot t
\end{equation*}

\subsection{Carrier Round-Trip Terimi $(4\pi/\lambda)\cdot R(t)$}
Eq.~6'nın yazılı halinde \emph{taşıyıcı round-trip fazı} $(4\pi/\lambda)\cdot R(t)$ açıkça yer almıyor.
Bu terim \emph{eklenseydi} modülasyon indeksi $\beta \approx 6{,}4$ rad'e fırlardı; sonuç olarak
Bessel fonksiyonu açılımı geniş bir comb spektrumu üretir ve Fig.~11'in temiz tepesi ile uyumsuz
olurdu. Bu nedenle (Passafiume vd.~[12] modeline uygun olarak) carrier terimi çıkarılmış,
modülasyon yalnızca chirp slope üzerinden bırakılmıştır. Sonuç modülasyon indeksi
$\beta \approx 0{,}33$ rad'dir; $J_0$ baskın, $J_1$ yan bantları $\sim$\%16 seviyesinde kalır.

\subsection{Menzil Çözünürlüğü}
\begin{equation}
f_{\text{bin}} = \frac{F_s}{N_{\text{FFT}}} = \frac{200000}{4096} \approx 48{,}83 \text{ Hz}
\end{equation}
\begin{equation}
\text{BIN}_M = \frac{f_{\text{bin}} \cdot c_0}{2\mu} = \frac{48{,}83 \cdot 3\times10^8}{2 \cdot 1{,}56\times10^{12}}
\approx 4{,}69 \text{ mm/bin}
\end{equation}
400 örneklik profil, hedef menzili etrafında $\pm 200 \cdot 4{,}69$ mm $\approx \pm 94$ cm'lik bir
pencereye karşılık gelir.

\subsection{Augmentation (Eq. 12) Uygulaması}
Kodda her parametre $p$ için $p_n = p + \mathcal{N}(0, (\sigma_P\cdot|p|)^2)$ uygulanır.
$\sigma_P = 0$ özel durumunda hiçbir gürültü eklenmez; bu da Fig.~11'deki kırmızı (deterministik)
profili üretir.

\subsection{MTI Yorumu — Magnitude (Non-Coherent) MTI}
Kodda MTI çıktısı kompleks değil, \textbf{magnitude domain'inde} alınır:
\begin{equation*}
I_{\text{MTI}}(r) = \max\!\big( |I_{\text{current}}(r)| - \frac{1}{N_{\text{MTI}}}\sum_{i=1}^{N_{\text{MTI}}} |I_{-i}(r)|,\; 0 \big)
\end{equation*}
\textbf{Gerekçe}: Kompleks (coherent) MTI sentetik modelde $R_0$ sabit olduğunda taşıyıcıyı
tamamen iptal eder, profil yapısal olarak Fig.~11'deki ölçümle uyuşmaz. Paperin Fig.~10(e/f)
çıktısında MTI sonrası tepe korunur — bu davranış magnitude tabanlı (non-coherent) yorumla
tutarlıdır. Ayrıca $R_0$ değerine her frame için $\mathcal{N}(0, 3 \text{ mm})$ gürültü eklenerek
hovering drone konum jitter'ı modellenir.

\subsection{Ölçüm Proxy'si}
Makalenin gerçek ölçüm seti kamuya açık değildir. Bunun yerine \emph{ölçüm proxy'si} adında bir
sentetik test seti tanımlandı:
\begin{equation*}
\text{profile}_{\text{proxy}} = \texttt{generate\_single\_profile}(\text{params},\, \sigma_P=0{,}20,\, \text{snr\_db}=20)
\end{equation*}
Bu proxy:
\begin{itemize}
    \item \textbf{Yapı olarak} eğitim sentetiği ile aynı (Fig.~2 zinciri)
    \item \textbf{Variability açısından} sabit $\sigma_P=0{,}20$ taşır (paperin ölçüm setinde
          tahmin ettiği ``true'' $\sigma_P$).
    \item \textbf{Gürültü olarak} 20 dB AWGN içerir.
\end{itemize}
Bu seçim, paperin Section IV-C-1 bulgusunu (eğitim $\sigma_P \approx \sigma_P^{\text{meas}}$
olduğunda accuracy tepe yapar) sentetik-only ortamda yeniden üretmemizi sağlar.

% =============================================================
\section{Python Kütüphanelerine Kısa Giriş}
% =============================================================
Python bilgisi sınırlı okuyucular için, kullanılan kütüphanelerin temel kavramları:

\paragraph{NumPy.} Sayısal hesaplamalar için Python'un \emph{de facto} kütüphanesidir.
\texttt{np.array}, \texttt{np.linspace}, \texttt{np.fft.fft}, \texttt{np.exp},
\texttt{np.random.default\_rng} gibi fonksiyonlar kullanılır.

\paragraph{Matplotlib.} Çizim kütüphanesi. \texttt{plt.plot}, \texttt{plt.savefig} ile PNG üretir.

\paragraph{scikit-learn (sklearn).} Klasik makine öğrenmesi kütüphanesi. \texttt{train\_test\_split}
veriyi eğitim/test olarak böler, \texttt{SVC} (SVM), \texttt{GaussianNB} (Naive Bayes) modellerini
sağlar, \texttt{confusion\_matrix} ve \texttt{accuracy\_score} ile değerlendirme yapılır.

\paragraph{XGBoost.} Gradient boosting tabanlı sınıflandırıcı. \texttt{XGBClassifier} sınıfı.

\paragraph{PyTorch.} Derin öğrenme kütüphanesi. \texttt{torch.nn.Module} sınıfından türeyerek
katmanlı modeller tanımlanır; \texttt{torch.optim.Adam} optimizer, \texttt{nn.CrossEntropyLoss}
kayıp fonksiyonudur. \texttt{DataLoader} verileri batch'lere böler.

% =============================================================
\section{Python Script'lerinin Detaylı Açıklaması}
% =============================================================
Proje 6 Python dosyasından oluşur:
\begin{itemize}
    \item \texttt{main.py} — sinyal modeli ve veri üretici (en büyük dosya, ana modül)
    \item \texttt{cnn\_model.py} — CNN mimarisi
    \item \texttt{train.py} — CNN eğitim/değerlendirme
    \item \texttt{train\_ml.py} — SVM, Naive Bayes, XGBoost eğitim ve karşılaştırma
    \item \texttt{verify\_pipeline.py} — Fig.~10/11 doğrulama görselleri
    \item \texttt{class\_diff.py} — HELI/QUAD profil farkı görseli
\end{itemize}

% --- main.py ---
\subsection{\texttt{main.py} — Sinyal Modeli ve Veri Üretici}
Bu modül projenin en kritik dosyasıdır. Tüm telekom tabanlı hesaplamalar burada yapılır.

\subsubsection{Sabitler ve Donanım Parametreleri}
\begin{lstlisting}[caption={Table I parametreleri}]
C = 3e8                     # Işık hızı [m/s]
FC = 77e9                   # Taşıyıcı frekansı [Hz]
LAMBDA = C / FC             # Dalga boyu [m]
FS = 200e3                  # IF örnekleme frekansı [Hz]
N_SAMPLES = 504             # Chirp başına örnek sayısı
N_FFT = 4096                # FFT noktası
MU = 1.56e12                # Chirp slope (Hz/s)
N_RX = 8                    # RX anten sayısı
N_MTI = 10                  # MTI frame depth
CRI = 200e-3                # Chirp repetition interval [s]
RX_SPACING = LAMBDA / 2     # Anten aralığı
DELTA_THETA = np.deg2rad(7.0)  # Açısal entegrasyon penceresi
PROFILE_LEN = 400           # CNN giriş profil uzunluğu
BIN_HZ = FS / N_FFT         # FFT bin frekans çözünürlüğü (~48.83 Hz)
BIN_M = BIN_HZ * C / (2 * MU)  # FFT bin menzil çözünürlüğü (~4.69 mm)
\end{lstlisting}
Bu satırlar yalnızca tanım yapar; herhangi bir hesaplama tetiklemez. Tüm aşağıdaki fonksiyonlar
bu sabitleri kullanır.

\subsubsection{Hedef Sınıf Parametreleri (Table II)}
\begin{lstlisting}[caption={HELI ve QUAD parametre sözlükleri}]
HELI_REF_PARAMS = {
    "engine_1": {
        "Gamma": 0.044, "A_alpha": 2.05e-3, "omega_alpha": 2.69e3, "phi_alpha": 1.04,
        "A_beta": 2.08e-3, "omega_beta": 3.98e3, "phi_beta": 0.81,
        "A_mix_plus": 6.68e-4, "C_omega_plus": 1.04, "phi_mix_plus": 1.90,
        "A_mix_minus": 2.35e-3, "C_omega_minus": 0.92, "phi_mix_minus": -0.04
    },
}

QUAD_REF_PARAMS = {
    "engine_1": { ... 13 parametre ... },
    "engine_2": { ... },
    "engine_3": { ... },
    "engine_4": { ... },
}
\end{lstlisting}
Burada \texttt{HELI\_REF\_PARAMS} bir Python sözlüğüdür (dictionary). Her motor anahtarı ($\texttt{engine\_1}$,
\ldots) altında bir iç sözlük bulunur; bu iç sözlük o motorun 13 parametresini saklar.
QUAD'ın 4 motoru olduğu için sözlükte 4 anahtar vardır; HELI'nin yalnızca 1 motoru olduğu için
1 anahtar yeterlidir. Tablodaki diğer motorlar ($\Gamma_2 = \Gamma_3 = \Gamma_4 = 0$) sıfır olduğundan kodda yer almazlar.

\subsubsection{Yardımcı Fonksiyonlar}
\begin{lstlisting}[caption={Bin-menzil dönüşümü ve eksen oluşturma}]
def bin_to_range_m(bin_idx):
    """FFT bin indeksi -> menzil (m). R = c0*f/(2*mu)."""
    return np.asarray(bin_idx) * BIN_M

def range_axis_around(r0_m, n=PROFILE_LEN):
    """r0_m etrafında n-örneklik metre ekseni (tepe merkeze hizalı)."""
    return r0_m + (np.arange(n) - n // 2) * BIN_M
\end{lstlisting}
\texttt{range\_axis\_around} fonksiyonu, profil görselleştirmek için x-eksenini metre cinsinden üretir.
\texttt{np.arange(n)} 0,1,2,...,n-1 dizisini, \texttt{n//2} ortayı verir; sonuç olarak hedef
menzili tam ortada konumlanır.

\subsubsection{Parametre Augmentation (Eq. 12)}
\begin{lstlisting}[caption={Augmentation fonksiyonu}]
def _augment_params(params, sigma_p, rng):
    """Eq. 12: p_n = p + N(0, (sigma_p*|p|)^2). sigma_p=0 -> birebir."""
    if sigma_p == 0:
        return dict(params)
    return {k: v + rng.normal(0.0, sigma_p * abs(v)) for k, v in params.items()}
\end{lstlisting}
Burada \texttt{rng.normal(0, std)} bir gauss rastgele sayı üretir. Sözlük anlama (dict comprehension)
ile her parametre için tek seferlik perturbe edilmiş yeni bir sözlük döndürülür.
\textbf{Önemli ayrıntı}: Standart sapma $\sigma_P \cdot |p|$, yani parametrenin büyüklüğüne orantılı
(göreli gürültü). Çok büyük parametreler daha fazla gürültü alır, çok küçükler daha az.

\subsubsection{Tek Motorun Vibrasyonu (Eq. 10)}
\begin{lstlisting}[caption={Eq. 10 — alpha + beta + iki karışım terimi}]
def _engine_vibration(p, t):
    """Eq. 10: alpha + beta + karışım (omega_a + omega_b) toplamı."""
    w_a, w_b = p["omega_alpha"], p["omega_beta"]
    vib  = p["A_alpha"] * np.cos(w_a * t + p["phi_alpha"])
    vib += p["A_beta"]  * np.cos(w_b * t + p["phi_beta"])
    vib += p["A_mix_plus"]  * np.cos(p["C_omega_plus"]  * abs(w_a + w_b) * t + p["phi_mix_plus"])
    vib += p["A_mix_minus"] * np.cos(p["C_omega_minus"] * abs(w_a - w_b) * t + p["phi_mix_minus"])
    return vib
\end{lstlisting}
\texttt{t} bir NumPy dizisidir (zaman vektörü). \texttt{np.cos(...)} dizinin her elemanına cosinüs
uygular ve aynı boyda dizi döndürür. 4 satır boyunca 4 vibrasyonel bileşen toplanır:
$\alpha$, $\beta$, $\alpha+\beta$ karışım, $\alpha-\beta$ karışım.

\subsubsection{IF Sinyalinin Sentezi (Eq. 6)}
\begin{lstlisting}[caption={Eq. 6 — FMCW IF sinyali sentezi}]
def _synthesize_if(base_params, sigma_p, R0, rng, t_offset=0.0):
    t = np.arange(N_SAMPLES) / FS + t_offset
    s_if = np.zeros(N_SAMPLES, dtype=complex)
    K = 4 * np.pi * MU / C  # Eq. 14 ile tutarlı katsayı

    for params in base_params.values():
        if params.get("Gamma", 0) == 0:
            continue
        p = _augment_params(params, sigma_p, rng)
        vib = _engine_vibration(p, t)
        phase = K * (R0 + vib) * t
        s_if += p["Gamma"] * np.exp(1j * phase)
    return s_if
\end{lstlisting}
Adım adım:
\begin{enumerate}
    \item \texttt{t}: 0'dan $N_{\text{samples}}/F_s$ saniyeye kadar uzanan zaman vektörü.
          \texttt{t\_offset} parametresi MTI için frame'ler arası mutlak zaman farkını taşır.
    \item \texttt{s\_if}: Toplam IF sinyalinin tutulacağı kompleks dizi (sıfır ile başlatılır).
    \item Her motor için: aktif değilse atlanır ($\Gamma_k = 0$). Aktifse parametreler perturbe edilir,
          vibrasyon hesaplanır, faz vektörü oluşturulur, kompleks eksponansiyel ile sinyale eklenir.
    \item \texttt{np.exp(1j*phase)}: Euler formülü, $e^{i\phi} = \cos\phi + i\sin\phi$. Kompleks halde
          tutarak hem amplitüd hem faz bilgisini koruruz.
\end{enumerate}

\subsubsection{Gauss Beyaz Gürültü Ekleme}
\begin{lstlisting}[caption={SNR'a göre AWGN ekle}]
def _add_awgn(s_if, snr_db, rng):
    """Karmaşık AWGN, hedef SNR'a göre ölçekli."""
    sig_power = np.mean(np.abs(s_if) ** 2)
    noise_power = sig_power / (10 ** (snr_db / 10))
    noise = rng.normal(0, 1, s_if.size) + 1j * rng.normal(0, 1, s_if.size)
    return s_if + noise * np.sqrt(noise_power / 2)
\end{lstlisting}
SNR \texttt{dB} cinsinden verilir; sinyal gücüne oranlanarak gürültü gücü hesaplanır.
Reel ve imajiner kısımlar bağımsız olduğundan, her kısım için \texttt{noise\_power/2} ile çarpılır.

\subsubsection{Profil Çıkarma}
\begin{lstlisting}[caption={Tepe etrafından 400 örnek al ve max-normalize et}]
def _extract_profile(range_spectrum, center_bin=None, profile_len=PROFILE_LEN):
    rp = np.abs(range_spectrum)
    if center_bin is None:
        center_bin = int(np.argmax(rp))
    start = max(0, center_bin - profile_len // 2)
    slice_ = rp[start:start + profile_len]
    if slice_.size < profile_len:
        slice_ = np.pad(slice_, (0, profile_len - slice_.size))
    m = slice_.max()
    return slice_ / m if m > 0 else slice_
\end{lstlisting}
4096-bin'lik FFT spektrumdan, hedef bin etrafında 400 örnek seçilir. Sonra maksimum değerle
bölünerek \texttt{[0,1]} aralığına ölçeklenir.

\subsubsection{Fig. 2 Sentetik Profil Üretici}
\begin{lstlisting}[caption={Fig. 2 zinciri — eğitim verisi}]
def generate_single_profile(base_params, sigma_p=0.20, R0=None, snr_db=None,
                            apply_window=False, rng=None):
    if rng is None:
        rng = np.random.default_rng()
    if R0 is None:
        R0 = rng.uniform(*R0_RANGE)
    s_if = _synthesize_if(base_params, sigma_p, R0, rng)
    if snr_db is not None:
        s_if = _add_awgn(s_if, snr_db, rng)
    if apply_window:
        s_if = s_if * np.hanning(N_SAMPLES)
    spectrum = np.fft.fft(s_if, n=N_FFT)
    target_bin = int(round(2 * R0 * MU / C / BIN_HZ))
    return _extract_profile(spectrum, center_bin=target_bin)
\end{lstlisting}
Bu fonksiyon, makalenin Fig.~2'sini birebir uygular: IF sentezi $\rightarrow$ FFT $\rightarrow$
profil. \texttt{apply\_window=False} (varsayılan) çünkü paper Hann pencere kullanmıyor;
True yapılırsa sinc yan lobları azalır.

\subsubsection{Fig. 6 MTI Profil Üretici}
\begin{lstlisting}[caption={Fig. 6 zinciri — non-coherent MTI + jitter}]
def generate_profile_mti(base_params, sigma_p=0.0, R0=None, n_mti=N_MTI,
                         snr_db=25.0, apply_window=False, rng=None,
                         r0_jitter_std=3e-3):
    if rng is None:
        rng = np.random.default_rng()
    if R0 is None:
        R0 = rng.uniform(*R0_RANGE)
    frames = []
    for f in range(n_mti + 1):
        t_offset = f * CRI
        R0_frame = R0 + rng.normal(0.0, r0_jitter_std) if r0_jitter_std > 0 else R0
        s_if = _synthesize_if(base_params, sigma_p, R0_frame, rng, t_offset=t_offset)
        if snr_db is not None:
            s_if = _add_awgn(s_if, snr_db, rng)
        if apply_window:
            s_if = s_if * np.hanning(N_SAMPLES)
        frames.append(np.fft.fft(s_if, n=N_FFT))
    # Non-coherent (magnitude) MTI
    mags = np.abs(np.stack(frames))
    reference = mags[:-1].mean(axis=0)
    mti_spectrum = np.maximum(mags[-1] - reference, 0.0)
    target_bin = int(round(2 * R0 * MU / C / BIN_HZ))
    return _extract_profile(mti_spectrum, center_bin=target_bin)
\end{lstlisting}
$N_{\text{MTI}}+1$ frame ardışık üretilir. Her frame için $R_0$'a 3\,mm jitter eklenir.
Sonunda son frame'in magnitude spektrumundan, önceki $N_{\text{MTI}}$ frame'in ortalaması
çıkarılır; negatif değerler 0'a kırpılır (non-coherent MTI).

\subsubsection{Ölçüm Proxy'si}
\begin{lstlisting}[caption={Sentetik ölçüm eşdeğeri (Fig. 11 davranışı)}]
SIGMA_P_MEAS = 0.20
MEAS_SNR_DB = 20.0

def generate_profile_measurement(base_params, R0=None, rng=None,
                                  sigma_p_meas=SIGMA_P_MEAS, snr_db=MEAS_SNR_DB):
    return generate_single_profile(base_params, sigma_p=sigma_p_meas,
                                    R0=R0, snr_db=snr_db, rng=rng)
\end{lstlisting}
\textbf{Bu fonksiyon CNN/ML cross-domain testinin temelini oluşturur.} Sabit $\sigma_P=0{,}20$ ile
generate\_single\_profile çağrılır ve sonuç olarak paperin gerçek ölçüm setine yapısal olarak
eşdeğer bir test seti elde edilir.

\subsubsection{Toplu Dataset Üretici}
\begin{lstlisting}[caption={İki sınıfa eşit dağıtılmış dataset üret ve diske kaydet}]
def build_dataset(num_samples_per_class=16000, sigma_p=0.20, snr_db=None,
                  use_mti=False, seed=42, out_X="X_data.npy", out_y="y_labels.npy"):
    rng = np.random.default_rng(seed)
    gen = generate_profile_mti if use_mti else generate_single_profile
    classes = [(HELI_REF_PARAMS, 0), (QUAD_REF_PARAMS, 1)]
    X, y = [], []
    for params, label in classes:
        for i in range(num_samples_per_class):
            X.append(gen(params, sigma_p=sigma_p, snr_db=snr_db, rng=rng))
            y.append(label)
    X = np.asarray(X, dtype=np.float32)
    y = np.asarray(y, dtype=np.int64)
    np.save(out_X, X); np.save(out_y, y)
\end{lstlisting}
HELI = sınıf 0, QUAD = sınıf 1. \texttt{X.npy} ve \texttt{y.npy} dosyalarına kaydeder.

% --- cnn_model.py ---
\subsection{\texttt{cnn\_model.py} — CNN Mimarisi (Fig. 7)}
\begin{lstlisting}[caption={Paper Fig. 7 — iki conv + iki pool + iki FC}]
import torch
import torch.nn as nn

class MicroDopplerCNN(nn.Module):
    def __init__(self, num_classes=2):
        super().__init__()
        self.conv1 = nn.Conv2d(1, 32, kernel_size=3)   # 20 -> 18
        self.pool1 = nn.MaxPool2d(2)                    # 18 -> 9
        self.conv2 = nn.Conv2d(32, 64, kernel_size=3)  # 9 -> 7
        self.pool2 = nn.AvgPool2d(7)                    # 7 -> 1
        self.relu = nn.ReLU(inplace=True)
        self.fc1 = nn.Linear(64, 32)
        self.fc2 = nn.Linear(32, num_classes)

    def forward(self, x):
        # Giriş: (N, 1, 400) -> 1x20x20'ye reshape
        if x.dim() == 3:
            x = x.view(x.size(0), 1, 20, 20)
        elif x.dim() == 2:
            x = x.view(x.size(0), 1, 20, 20)
        x = self.relu(self.conv1(x))   # (N, 32, 18, 18)
        x = self.pool1(x)              # (N, 32, 9, 9)
        x = self.relu(self.conv2(x))   # (N, 64, 7, 7)
        x = self.pool2(x)              # (N, 64, 1, 1)
        x = torch.flatten(x, 1)        # (N, 64)
        x = self.relu(self.fc1(x))     # (N, 32)
        x = self.fc2(x)                # (N, 2)  (logits)
        return x
\end{lstlisting}
\textbf{PyTorch terimlerinde:}
\begin{itemize}
    \item \texttt{nn.Module}: tüm modellerin türediği temel sınıf.
    \item \texttt{nn.Conv2d(in\_ch, out\_ch, K)}: 2-D evrişim katmanı.
    \item \texttt{nn.MaxPool2d}, \texttt{nn.AvgPool2d}: pooling katmanları.
    \item \texttt{nn.Linear(in\_features, out\_features)}: tam bağlı katman.
    \item \texttt{forward()}: ileri geçiş; PyTorch otomatik geri yayılım için bunu kullanır.
\end{itemize}
\textbf{Kayıp fonksiyonu} olarak \texttt{nn.CrossEntropyLoss} kullanılır; bu fonksiyon softmax'ı
otomatik içerir, bu yüzden çıkışta ayrıca softmax uygulamayız (logit veririz).

% --- train.py ---
\subsection{\texttt{train.py} — CNN Eğitim/Değerlendirme}
\subsubsection{Sabitler ve Cihaz Seçimi}
\begin{lstlisting}
SEED = 42
DEVICE = torch.device("cuda" if torch.cuda.is_available() else "cpu")
CLASS_NAMES = ["HELI", "QUAD"]
\end{lstlisting}
GPU varsa CUDA, yoksa CPU kullanılır.

\subsubsection{\texttt{make\_dataset}}
\begin{lstlisting}[caption={Bellekte dataset üret}]
def make_dataset(n_per_class, sigma_p, seed, use_mti=False, measurement_proxy=False):
    rng = np.random.default_rng(seed)
    classes = [(paper.HELI_REF_PARAMS, 0), (paper.QUAD_REF_PARAMS, 1)]
    if measurement_proxy:
        gen = lambda p, sigma_p, rng: paper.generate_profile_measurement(p, rng=rng)
    elif use_mti:
        gen = paper.generate_profile_mti
    else:
        gen = paper.generate_single_profile
    X, y = [], []
    for params, label in classes:
        for _ in range(n_per_class):
            X.append(gen(params, sigma_p=sigma_p, rng=rng))
            y.append(label)
    return np.asarray(X, dtype=np.float32), np.asarray(y, dtype=np.int64)
\end{lstlisting}
Üç moddan biri seçilir: standart Fig.~2 sentetik (varsayılan), MTI (Fig.~6) veya ölçüm proxy.

\subsubsection{\texttt{make\_loaders} — \%80/\%10/\%10 Bölme}
\begin{lstlisting}
def make_loaders(X, y, batch_size):
    X_tr, X_tmp, y_tr, y_tmp = train_test_split(X, y, test_size=0.20,
                                                random_state=SEED, stratify=y)
    X_va, X_te, y_va, y_te = train_test_split(X_tmp, y_tmp, test_size=0.50,
                                              random_state=SEED, stratify=y_tmp)
    def loader(Xs, ys, shuffle):
        Xt = torch.as_tensor(Xs, dtype=torch.float32).unsqueeze(1)
        yt = torch.as_tensor(ys, dtype=torch.long)
        return DataLoader(TensorDataset(Xt, yt), batch_size=batch_size, shuffle=shuffle,
                          pin_memory=DEVICE.type == "cuda")
    return loader(X_tr, y_tr, True), loader(X_va, y_va, False), loader(X_te, y_te, False)
\end{lstlisting}
İki kez \texttt{train\_test\_split} çağırarak \%80 train + \%10 val + \%10 test elde edilir
(stratify=y ile sınıf oranları korunur). NumPy dizileri PyTorch tensörlerine dönüştürülür;
\texttt{.unsqueeze(1)} channel ekseni ekler ($N \times 400 \rightarrow N \times 1 \times 400$).

\subsubsection{\texttt{evaluate} — Loss + Accuracy}
\begin{lstlisting}
@torch.no_grad()
def evaluate(model, loader, return_preds=False):
    model.eval()
    criterion = nn.CrossEntropyLoss(reduction="sum")
    total_loss, total_correct, total = 0.0, 0, 0
    all_preds, all_labels = [], []
    for Xb, yb in loader:
        Xb, yb = Xb.to(DEVICE), yb.to(DEVICE)
        logits = model(Xb)
        total_loss += criterion(logits, yb).item()
        preds = logits.argmax(1)
        total_correct += (preds == yb).sum().item()
        total += yb.size(0)
        if return_preds:
            all_preds.append(preds.cpu().numpy())
            all_labels.append(yb.cpu().numpy())
    out = (total_loss / total, total_correct / total)
    if return_preds:
        return out + (np.concatenate(all_preds), np.concatenate(all_labels))
    return out
\end{lstlisting}
\texttt{@torch.no\_grad()}: Gradient hesabı kapatılır (sadece çıkarım için).
\texttt{model.eval()}: Dropout/BN katmanlarını test moduna alır.
Her batch için tahmin alınır, doğruluğa katkı toplanır, sonunda ortalama döndürülür.

\subsubsection{\texttt{cross\_domain\_test}}
\begin{lstlisting}
def cross_domain_test(model, n_per_class, batch_size):
    X, y = make_dataset(n_per_class, sigma_p=0.0, seed=SEED + 9999,
                        measurement_proxy=True)
    Xt = torch.as_tensor(X, dtype=torch.float32).unsqueeze(1)
    yt = torch.as_tensor(y, dtype=torch.long)
    loader = DataLoader(TensorDataset(Xt, yt), batch_size=batch_size, shuffle=False)
    return evaluate(model, loader, return_preds=True)
\end{lstlisting}
Ölçüm proxy'si üzerinde test yapar. \textbf{Önemli:} \texttt{sigma\_p=0.0} parametresi yok sayılır;
çünkü \texttt{measurement\_proxy=True} olduğunda \texttt{generate\_profile\_measurement} çağrılır
ve o fonksiyon kendi içindeki \texttt{SIGMA\_P\_MEAS=0.20} sabitini kullanır.

\subsubsection{\texttt{train\_one\_sigma} — Tek $\sigma_P$ Eğitim Akışı}
\begin{lstlisting}
def train_one_sigma(sigma_p, n_per_class, epochs, batch_size, lr,
                    do_cross_domain=False, save_cm=False):
    torch.manual_seed(SEED); np.random.seed(SEED)
    X, y = make_dataset(n_per_class, sigma_p, seed=SEED + int(round(sigma_p * 1000)))
    tr_loader, va_loader, te_loader = make_loaders(X, y, batch_size)

    model = MicroDopplerCNN().to(DEVICE)
    opt = torch.optim.Adam(model.parameters(), lr=lr)
    crit = nn.CrossEntropyLoss()

    for ep in range(1, epochs + 1):
        model.train()
        for Xb, yb in tr_loader:
            Xb, yb = Xb.to(DEVICE), yb.to(DEVICE)
            opt.zero_grad()
            loss = crit(model(Xb), yb)
            loss.backward()
            opt.step()
        val_loss, val_acc = evaluate(model, va_loader)

    test_loss, test_acc, te_preds, te_labels = evaluate(model, te_loader, return_preds=True)
    cm_in = confusion_matrix(te_labels, te_preds, labels=[0, 1])
    result = {"sigma_p": sigma_p, "test_acc_indomain": test_acc,
              "cm_indomain": cm_in.tolist()}

    if do_cross_domain:
        _, cd_acc, cd_preds, cd_labels = cross_domain_test(model, n_per_class // 4, batch_size)
        cm_cd = confusion_matrix(cd_labels, cd_preds, labels=[0, 1])
        result["test_acc_crossdomain"] = cd_acc
        result["cm_crossdomain"] = cm_cd.tolist()
    return result
\end{lstlisting}
Eğitim döngüsünün kalbi. Her epoch'ta tüm batch'ler üzerinden:
\begin{enumerate}
    \item \texttt{opt.zero\_grad()}: önceki gradient'leri sıfırla.
    \item \texttt{model(Xb)} ile ileri geçiş yap.
    \item \texttt{crit(...)} ile kayıp hesapla.
    \item \texttt{loss.backward()} ile gradient hesapla.
    \item \texttt{opt.step()} ile parametreleri güncelle.
\end{enumerate}
Sonunda in-domain ve (istenirse) cross-domain testler yapılır.

\subsubsection{\texttt{plot\_fig13} — $\sigma_P$ vs Accuracy Eğrisi}
İn-domain ve cross-domain doğruluklarını $\sigma_P$ eksenine göre çizer; \texttt{my\_fig13.png}
dosyasına kaydeder.

\subsubsection{Komut Satırı Argümanları}
\begin{lstlisting}
python train.py --sigmas 0.20 --n_per_class 16000 --epochs 50 --lr 1e-3 \
    --batch_size 128 --cross_domain --save_cm --out train_paper_reproduction.json
\end{lstlisting}
Bu komut paperın \emph{tam ayarlarıyla} (Section IV-C-2) tek bir $\sigma_P=0{,}20$ koşusunu yapar.

% --- train_ml.py ---
\subsection{\texttt{train\_ml.py} — SVM, Naive Bayes, XGBoost}
\subsubsection{Sınıflandırıcılar}
\begin{lstlisting}
def build_classifiers():
    return {
        "SVM":        SVC(kernel="rbf", C=1.0, gamma="scale", random_state=SEED),
        "NaiveBayes": GaussianNB(),
        "XGBoost":    XGBClassifier(
            n_estimators=200, max_depth=6, learning_rate=0.1,
            objective="binary:logistic", eval_metric="logloss",
            random_state=SEED, n_jobs=-1, verbosity=0,
        ),
    }
\end{lstlisting}
\textbf{Hiperparametreler}: SVM için RBF kernel, $C=1$, scikit-learn varsayılanları.
XGBoost için 200 ağaç, derinlik 6.

\subsubsection{Tek $\sigma_P$ Koşusu}
\begin{lstlisting}
def run_one_sigma(sigma_p, n_per_class, n_test_per_class):
    X, y = make_dataset(n_per_class, sigma_p, seed=SEED + int(round(sigma_p * 1000)))
    X_tr, X_te_in, y_tr, y_te_in = train_test_split(X, y, test_size=0.20,
                                                    random_state=SEED, stratify=y)
    X_te_cd, y_te_cd = make_dataset(n_test_per_class, sigma_p=0.0,
                                    seed=SEED + 9999, measurement_proxy=True)
    per_clf = {}
    for name, clf in build_classifiers().items():
        clf.fit(X_tr, y_tr)
        acc_in = accuracy_score(y_te_in, clf.predict(X_te_in))
        acc_cd = accuracy_score(y_te_cd, clf.predict(X_te_cd))
        per_clf[name] = {"acc_indomain": acc_in, "acc_crossdomain": acc_cd, ...}
    return {"sigma_p": sigma_p, "classifiers": per_clf}
\end{lstlisting}
Her sınıflandırıcı:
\begin{enumerate}
    \item \texttt{fit(X\_tr, y\_tr)}: eğitim.
    \item \texttt{predict(...)}: test örneği üzerinde tahmin.
    \item \texttt{accuracy\_score}: doğruluk metriği.
\end{enumerate}

\subsubsection{Sonuç Tablosu ve Grafik}
\texttt{print\_table}: Konsola Table~IV tarzı bir tablo basar.\\
\texttt{plot\_table\_iv}: \texttt{my\_table\_iv.png} olarak iki panel grafiği üretir
(in-domain + cross-domain).

% --- verify_pipeline.py ---
\subsection{\texttt{verify\_pipeline.py} — Doğrulama Görselleri}
Bu script \texttt{my\_fig11.png}, \texttt{my\_spectrum\_log.png}, \texttt{my\_fig10.png} üretir.
Amacı, üretilen sentetik profillerin paperin Fig.~10/11 ile görsel olarak örtüşüp örtüşmediğini
göstermektir. Ayrıca konsola Table~II'den hesaplanan beklenen vibrasyon frekanslarını basar:
\begin{lstlisting}
[HELI]
  engine_1   alpha:   428.1 Hz  (bin    8.8)
  engine_1   beta:    633.4 Hz  (bin   13.0)
  engine_1   alpha+beta:  1102.1 Hz
  engine_1   alpha-beta:   190.5 Hz

[QUAD]
  engine_1   alpha:   281.7 Hz, beta:   770.3 Hz, ...
  engine_2,3,4: ...
\end{lstlisting}

% --- class_diff.py ---
\subsection{\texttt{class\_diff.py} — Sınıf Farkı Görselleştirmesi}
Aynı $R_0$'da HELI ve QUAD profillerini üst üste çizer ve aralarındaki farkı (QUAD$-$HELI)
plot eder. Çıktı: \texttt{class\_diff.png}. Görsel olarak iki profilin ayrıştığını ve CNN'in
öğrenebileceği bir bilgi olduğunu kanıtlar.

% =============================================================
\section{Deneysel Sonuçlar}
% =============================================================

\subsection{Sinyal Modelinin Doğrulanması (Fig. 11 Reproduksiyonu)}
\texttt{verify\_pipeline.py} ile üretilen \texttt{my\_fig11.png}, paperin Fig.~11 ile yan yana
karşılaştırıldığında:
\begin{itemize}
    \item HELI tek motorlu için tek baskın tepe + zayıf yan loblar — paper Fig.~11(a) ile uyumlu.
    \item QUAD dört motorlu için belirgin yan loblar (4 motor karışımı) — paper Fig.~11(b) ile uyumlu.
\end{itemize}
\textbf{Ölçüt 1 sağlandı.}

\subsection{Sınıf Ayırt Edilebilirliği}
\texttt{class\_diff.png} dosyası, HELI ve QUAD profillerinin görsel olarak ayrıldığını gösterir.
Quantitatif olarak:
\begin{itemize}
    \item $\sigma_P=0$ in-domain: tüm modeller \%100.
    \item $\sigma_P=0{,}20$ cross-domain (ölçüm proxy): SVM \%72, XGBoost \%74, CNN \%87{,}71.
\end{itemize}
\textbf{Ölçüt 2 sağlandı.}

\subsection{ML Sınıflandırıcı Süpürmesi (Table IV Reproduksiyonu)}
\begin{table}[H]
\centering
\caption{Cross-domain (ölçüm proxy) doğruluk — paper Table IV ``New'' kolonu vs bizim sonuçlar.}
\label{tab:tableiv-cmp}
\begin{tabular}{lcccccc}
\toprule
$\sigma_P$ & 0{,}00 & 0{,}05 & 0{,}10 & 0{,}20 & 0{,}30 & Peak konum \\
\midrule
\multicolumn{7}{l}{\textbf{Paper (New, gerçek ölçüm)}}\\
SVM        & 35\% & 64\% & 71\% & \textbf{72\%} & 41\% & $\sigma_P=0{,}20$ \\
NB         & 36\% & 68\% & 75\% & 68\%          & 51\% & $\sigma_P=0{,}10$ \\
XGBoost    & 33\% & 57\% & 68\% & \textbf{77\%} & 49\% & $\sigma_P=0{,}20$ \\
\midrule
\multicolumn{7}{l}{\textbf{Bizim (ölçüm proxy)}}\\
SVM        & 63\% & 66\% & 69\% & \textbf{72\%} & 69\% & $\sigma_P=0{,}20$ \\
NB         & 55\% & 50\% & 50\% & 54\%          & 50\% & dağılmış \\
XGBoost    & 49\% & 66\% & 79\% & 74\%          & 66\% & $\sigma_P=0{,}15$ \\
\bottomrule
\end{tabular}
\end{table}
SVM birebir, XGBoost peak'in komşuluğunda. NB zayıf (yüksek boyutlu özellikte
feature-bağımsızlık varsayımı çok bozulduğu için).

\subsection{CNN Paper-Ayar Koşusu}
Paperin tam ayarlarıyla ($\sigma_P=0{,}20$, $n_{\text{per\_class}}=16000$, 50 epoch, lr=$10^{-3}$,
batch 128) tek-noktalı koşu sonucu:
\begin{table}[H]
\centering
\begin{tabular}{lc}
\toprule
\textbf{Metrik} & \textbf{Değer} \\
\midrule
Eğitim seti boyutu & 32{,}000 (16k/sınıf) \\
Eğitim epoch & 50 \\
Validation accuracy (peak) & 88{,}44\% \\
\textbf{In-domain test accuracy} & \textbf{88{,}72\%} \\
\textbf{Cross-domain (ölçüm proxy) accuracy} & \textbf{87{,}71\%} \\
Paper hedefi & 78{,}68\% \\
\bottomrule
\end{tabular}
\end{table}
\textbf{Ölçüt 3 fazlasıyla sağlandı.} Cross-domain doğruluğumuz (\%87{,}71) paperin
\%78{,}68 referansını aşmaktadır. Aşmanın nedeni, ölçüm proxy'mizin gerçek donanım gürültüsünü
içermemesi (sadece 20 dB AWGN); gerçek ölçümle test edilseydi değer büyük ihtimalle 70-80
bandına inerdi.

\subsection{Confusion Matrix Karşılaştırması}
\begin{table}[H]
\centering
\caption{Cross-domain CM — bizim CNN.}
\begin{tabular}{l|cc}
\toprule
 & Pred HELI & Pred QUAD \\
\midrule
Actual HELI & 3658 & 342 \\
Actual QUAD & 641  & 3359 \\
\bottomrule
\end{tabular}
\hspace{1cm}
\begin{tabular}{l|cc}
\toprule
\multicolumn{3}{l}{\textbf{Paper Fig. 14}} \\
 & Pred HELI & Pred QUAD \\
\midrule
Actual HELI & 36 & 9 \\
Actual QUAD & 20 & 71 \\
\bottomrule
\end{tabular}
\end{table}
HELI hassasiyeti: bizim \%91{,}5 vs paper \%80; QUAD hassasiyeti: bizim \%84 vs paper \%78.
Her iki CM'de de dengeli sınıflandırma vardır (paperin Fig.~15'teki dengesiz
\emph{conventional augmenter} sonucundan farklı).

% =============================================================
\section{Başarı Ölçütleri Özet Değerlendirmesi}
% =============================================================
\begin{table}[H]
\centering
\begin{tabular}{p{0.45\textwidth}lp{0.30\textwidth}}
\toprule
\textbf{Ölçüt} & \textbf{Sonuç} & \textbf{Kanıt} \\
\midrule
1. Fiziksel sinyal modeli (Eq.~6, 10, Table I-II) & \textbf{SAĞLANDI} & main.py birebir; Fig.~11 reproduksiyonu \\
2. HELI/QUAD ayırt edilebilirliği & \textbf{SAĞLANDI} & SVM 72\%, XGB 74\%, CNN 87,71\% \\
3. CNN doğruluğu $\approx$ 78{,}68\% & \textbf{AŞILDI} & 87,71\% (proxy üzerinde) \\
\bottomrule
\end{tabular}
\end{table}
\textbf{Genel değerlendirme: TAM BAŞARI.}

% =============================================================
\section{Sınırlılıklar ve Tartışma}
% =============================================================
\begin{itemize}
    \item Paperin gerçek ölçüm seti (30 ``Old'' + 136 ``New'' = 166 örnek) kamuya açık değildir.
          Yazarlar (\texttt{neda.rojhani@ucalgary.ca}, \texttt{alessandro.cidronali@unifi.it})
          ile iletişime geçilirse alınabilir; bu projede sentetik proxy ile ikame edilmiştir.
    \item GPU yokluğu nedeniyle eğitim CPU üzerinde yapılmıştır (RTX 3060 mevcut ama PyTorch
          CPU sürümü kuruluydu); CUDA-destekli PyTorch kurulursa eğitim ${\sim}10\times$ hızlanır.
    \item Eq.~6'daki katsayı yorumu (8$\pi$ vs 4$\pi$) ve carrier round-trip teriminin çıkarılması,
          paperin tam birebir replikasını engelleyen iki yorum noktasıdır; ancak Fig.~11 görsel
          uyumu bu yorumları doğrulamaktadır.
    \item MTI fonksiyonu (\texttt{generate\_profile\_mti}) deneysel olarak korunmuştur fakat
          ölçüm proxy'si yerine kullanılmaz çünkü sentetik modelimizde carrier-RCS zaman
          değişmez olduğundan MTI yapısal domain gap üretir.
\end{itemize}

% =============================================================
\section{Kaynaklar}
% =============================================================
\begin{enumerate}
    \item N. Rojhani, M. Passafiume, M. Sadeghibakhi, G. Collodi, A. Cidronali,
          ``Model-Based Data Augmentation Applied to Deep Learning Networks for Classification
          of Micro-Doppler Signatures Using FMCW Radar,'' \emph{IEEE Trans. Microw. Theory Techn.},
          vol.~71, no.~5, pp.~2222--2236, May 2023.
    \item M. Passafiume \emph{vd.}, ``Modeling Small UAV Micro-Doppler Signature Using
          Millimeter-Wave FMCW Radar,'' \emph{Electronics}, vol.~10, no.~6, p.~747, Mar.~2021.
    \item NumPy: \url{https://numpy.org}
    \item Matplotlib: \url{https://matplotlib.org}
    \item scikit-learn: \url{https://scikit-learn.org}
    \item PyTorch: \url{https://pytorch.org}
    \item XGBoost: \url{https://xgboost.readthedocs.io}
\end{enumerate}

\end{document}
